% Cote 119: offprints, reports, courses, CV and letters (typed and printed part): archive cote 119, site page 242 (PDF page 243).
% typescript. Draft compared: none.
% AI-transcribed by Claude Opus 5.5 (claude-opus-5-5, adaptive thinking, effort high), Message Batches API, batch r119_full (msgbatch_01DwZNm5Y7YsKBWyAN4sFs8w), 2026-11-13; not yet checked by eye.
\scanpage{119/242}
\mspage{2}

changé, non seulement en France, mais bien entendu partout dans le monde.
Tout d'abord, avec les \uline{motivations} de l'enseignement, qui deviennent de
plus en plus utilitaires, ses \uline{besoins} ont également changé. Parmi les
étudiants, de plus en plus nombreux sont ceux qui n'ont pas besoin d'une
connaissance théorique plus ou moins exhaustive, allant jusqu'aux limites
du connu, de telle ou telle science, la Mathématique disons ; mais plutôt
de telle ou telle partie de la Mathématique, qu'ils auront à utiliser seu-%
lement à titre auxiliaire dans une variété de métiers techniques ou scien-%
tifiques possibles : ingénieur, physicien, chimiste, biologiste, économiste...
Ce sont donc de\xout{s} futurs utilisateurs de Mathématiques, non de futurs mathé-%
maticiens. Il leur est nécessaire de comprendre un formalisme mathématique
limité, en vue de pouvoir l'appliquer aux problèmes spécifiques de leurs
futurs métiers. Sous ces conditions, les questions de rigueur logique par
exemple (qui pour la Mathématique en tant que science pure sont et sans
doute resteront une exigence essentielle) deviennent accessoires, l'exigence
principale étant pour eux de savoir comprendre et surtout, appliquer correc-%
tement, les algorithmes qu'ils seront amenés à utiliser.

Pour répondre à cette plus grande variété de besoins, il convient
donc de dispenser une plus grande variété d'enseignements,- \uline{variété non seu-}%
\uline{lement par le nombre des théories mathématiques enseignées}, \uline{mais par l'optique}
\uline{dans laquelle elles sont enseignées}, qui doit être fonction des besoins
véritables des étudiants, et ne pas rester soumise à des d\unsure{è|é}sid\unsure{è|é}rata (de
rigueur, d'exhaustivité ou d'universalité) qui étaient justifiés il y a
cent ans et ne le sont plus aujourd'hui.

Parallèlement à cette différentiation des besoins des \uline{enseignés},
s'accentue nécessairement, qu'on le veuille ou non, la différenti ation
entre les deux rôles dévolus au \uline{maître}. Sur le plan qualitatif, c'est
surtout vrai pour l'enseignement à assurer à de futurs utilisateurs de la
Mathématique comme plus haut, qui n'auront rien à gagner (et qui risquent
fort d'y perdre) que cet enseignement soit assuré par des savants éminents,
ou simplement des chercheurs de valeur. En règle générale et sauf de rares
exceptions, plus grande sera la valeur d'un chercheur, c'est-à-dire plus

\margin{242}
